\documentclass{article}
\usepackage[utf8]{inputenc}
\usepackage[margin=0.75in]{geometry}

\begin{document}
	\begin{center}
    
    	% MAKE SURE YOU TAKE OUT THE SQUARE BRACKETS
    
		\LARGE{\textbf{Capstone Proposal}} \\
        \vspace{1em}
        \Large{Compressing LLMs with Holograms} \\
        \vspace{1em}
        \normalsize\textbf{Karapet Khachatryan} \\
        \normalsize{karapet\_khachatryan\@edu.aua.am} \\ % add hyperlink
        \vspace{1em}
        \normalsize{Supervisor: Suren Khachatryan} \\
        \vspace{1em}
        \normalsize{American University of Armenia, Yerevan, Armenia} \\
        \normalsize{Bachelors of Computer Science}
    
	\end{center}

    \begin{normalsize}
    
    	\section{Abstract}
        This capstone project investigates whether principles from optical holography can be adapted as an experimental compression method for large language model weight matrices. In classical holography, a three-dimensional object is not stored as a direct visual image, but as an interference pattern that contains information about the object wavefront, including both amplitude and phase. Inspired by this idea, the project treats neural network parameters as structured spatial data: weight matrices are converted into two-dimensional or three-dimensional geometric representations, encoded into a hologram-like image, and then decoded back into reconstructed matrices or point-cloud objects. The main research question is whether a compact two-dimensional holographic representation can preserve enough information to reconstruct model weights with acceptable numerical similarity.

        The proposed pipeline explores matrix normalization, matrix-to-surface mapping, VTP point-cloud generation, Gabor-style hologram encoding, amplitude and phase storage, and reconstruction quality measurement using metrics such as RMSE, MAE, relative error, cosine similarity, PSNR, and custom similarity scores. Unlike standard LLM compression methods such as pruning, quantization, and low-rank adaptation, this work does not begin by removing or simplifying parameters directly. Instead, it studies whether model weights can be re-represented through a physically inspired encoding process. The project is exploratory, and its value lies in testing the limits, weaknesses, and possible future potential of holographic representations for neural model compression.

        \section{Introduction}
        Large language models have become central to modern artificial intelligence, but their practical use is limited by the size of their parameter sets, memory requirements, and computational cost. Transformer-based models contain many large weight matrices, often repeated across attention and feed-forward layers. These matrices are usually stored directly as numerical arrays, which makes deployment expensive when the model has billions of parameters. For this reason, model compression has become an important research direction. Existing compression methods usually aim to reduce the number of stored bits, remove less important weights, approximate matrices with smaller structures, or transfer knowledge from a larger model to a smaller one.

        This capstone approaches the problem from a different perspective. Instead of asking only how many parameters can be removed or how many bits each weight can use, it asks whether neural weights can be transformed into another representational domain. The main inspiration is holography. Dennis Gabor introduced holography in 1948 as a method for recording and reconstructing wavefront information rather than simply capturing an ordinary intensity image. In holography, the recorded pattern may look visually meaningless, but it can reconstruct the appearance of a three-dimensional object when interpreted through the proper optical or numerical process. This is important for the project because it suggests a broader idea: information does not always need to be stored in the same form in which it is later used. A complex object can be encoded into an indirect representation and reconstructed afterward.

        The project applies this principle to model weights. A weight matrix can be interpreted as a surface, where the row and column indices define spatial coordinates and the weight value defines height, amplitude, depth, or phase. Multiple matrices can also be packed together into a larger spatial representation, or transformed onto curved surfaces such as hemispheres. After this conversion, the resulting object can be encoded into a hologram-like two-dimensional image. Decoding then attempts to reconstruct the original matrix or point cloud from the holographic representation.

        The research is not claiming that holographic encoding is already better than standard compression. Rather, it studies whether such an approach can preserve meaningful numerical structure while reducing storage size. Early experiments show that reconstruction can be very sensitive to hologram resolution, depth handling, angle selection, normalization, and whether only amplitude or both amplitude and phase are stored. A zero matrix reconstructs easily because it contains almost no structural variation, while more complex matrices expose the weakness of the current encoding and decoding process. This makes the project valuable as a research exploration: even failed or low-similarity reconstructions help identify what kind of information the holographic pipeline loses.

        The main goal of the capstone is therefore to design, implement, and evaluate a hologram-inspired encoding-decoding pipeline for neural weight matrices. The project uses simple geometric objects such as cubes and hemispheres as controlled test cases, then applies similar methods to language model weight matrices. The success of the method is evaluated through reconstruction metrics, visual inspection, file-size comparison, and analysis of whether depth, phase, and matrix structure are preserved. The broader contribution of the work is to connect ideas from Fourier optics, digital holography, and neural model compression into one experimental framework.

	   	\section{Literature Review}
    	\section{Problem Formulation}
        \section{Proposed Methodology}
        \section{Expected Results and Evaluation}
        \section{Timeline}
        \section{References}
    \end{normalsize}

  
\end{document}